\section{Algoritmos STL comunes}

\begin{lstlisting}[style=cpp]
void common_algorithms(vector<int> values) {
    sort(values.begin(), values.end());
    sort(values.rbegin(), values.rend());
    stable_sort(values.begin(), values.end()); // conserva el orden relativo de elementos que empatan (noo mas rapido si mas memo)

    int minimum = *min_element(values.begin(), values.end());
    int maximum = *max_element(values.begin(), values.end());
    auto [minimum_it, maximum_it] = minmax_element(values.begin(), values.end());

    ll sum = accumulate(values.begin(), values.end(), 0LL);
    int amount = count(values.begin(), values.end(), 5); // Cuantas veces aparece 5
    int positives = count_if(values.begin(), values.end(),
                             [](int value) { return value > 0; });
    bool all_positive = all_of(values.begin(), values.end(),
                               [](int value) { return value > 0; });

    reverse(values.begin(), values.end());
    rotate(values.begin(), values.begin() + values.size() / 2, values.end()); // {1, 2, 3, 4, 5, 6} -> {4, 5, 6, 1, 2, 3}
    (void)minimum; (void)maximum; (void)minimum_it; (void)maximum_it;
    (void)sum; (void)amount; (void)positives; (void)all_positive; //: no la uso; es intencional
}
\end{lstlisting}

\subsection{Bounds, duplicados y selecci\'on}

\begin{lstlisting}[style=cpp]
void ordered_range_algorithms(vector<int> values, int target) {
    sort(values.begin(), values.end());
    auto first_equal = lower_bound(values.begin(), values.end(), target);
    auto after_equal = upper_bound(values.begin(), values.end(), target);
    int occurrences = after_equal - first_equal;
    bool exists = binary_search(values.begin(), values.end(), target);

    values.erase(unique(values.begin(), values.end()), values.end());
    if (!values.empty()) {
        int middle = values.size() / 2;
        nth_element(values.begin(), values.begin() + middle, values.end());
        int median = values[middle];
        (void)median;
    }
    (void)occurrences; (void)exists;
}
\end{lstlisting}

\subsection{Permutaciones y sumas acumuladas}

\begin{lstlisting}[style=cpp]
void sequence_algorithms() {
    vector<int> permutation{1, 2, 3};
    do {
        // procesar permutation
    } while (next_permutation(permutation.begin(), permutation.end()));

    // partial_sum acumula con el tipo de entrada, no con el de prefix.
    // Usar ll; cada suma parcial debe entrar en ese tipo.
    vector<ll> values(5);
    iota(values.begin(), values.end(), 10LL); // 10,11,12,13,14
    vector<ll> prefix(values.size());
    partial_sum(values.begin(), values.end(), prefix.begin());
}
\end{lstlisting}

\warning{\texttt{lower\_bound}, \texttt{upper\_bound} y
\texttt{binary\_search} requieren el rango ordenado con el mismo comparador.}
